Производительность криптографического ядра зависит от каждого этапа обработки данных: от того, как лимбы ложатся в регистры процессора, до способа передачи буферов конечному пользователю. В этой главе разобраны технические решения, позволившие реализовать X25519 и Ed25519 на платформе .NET с максимальной эффективностью. Основное внимание уделено векторизации вычислений через AVX2, минимизации нагрузки на сборщик мусора и соблюдению принципов константного времени выполнения.
\subsection{Требования к программному комплексу}

На этапе проектирования были сформированы строгие функциональные и нефункциональные требования к разрабатываемой библиотеке, обусловленные спецификой современных высоконагруженных систем.

Функциональные требования:
\begin{enumerate}
    \item Реализация базовых вычислительных примитивов (скалярное умножение точки) на кривой Монтгомери (X25519) и искривленной кривой Эдвардса (Ed25519) в соответствии со спецификациями RFC 7748 \cite{RFC7748} и RFC 8032 \cite{RFC8032}.
    \item Реализация алгоритмов кодирования и декодирования (сериализации) точек эллиптических кривых для обеспечения совместимости со стандартами.
    \item Обеспечение параллельного вычисления пакетов из четырех независимых скалярных умножений за один вызов метода для оптимизации пропускной способности при массовых операциях.
\end{enumerate}

Нефункциональные требования (архитектурные атрибуты):
\begin{enumerate}
    \item Кроссплатформенность и управляемость (Managed Code): Библиотека должна быть написана исключительно на языке C\# без использования платформозависимых вызовов (DllImport / P/Invoke) к нативным библиотекам (C/C++), что гарантирует безопасность типов и простоту развертывания.
    \item Минимизация аллокаций (Zero Allocation): Критические участки кода, ответственные за выполнение криптографических преобразований, реализованы без выделения памяти в управляемой куче (Managed Heap). Это позволяет полностью исключить нагрузку на сборщик мусора (Garbage Collector) в процессе работы.   
    \item Константное время выполнения (Constant-time): Архитектура вычислительного ядра должна математически исключать ветвления, зависящие от секретных данных, для предотвращения атак по сторонним каналам (Side-Channel Attacks).
\end{enumerate}

\subsection{Многоуровневая архитектура библиотеки}

Для обеспечения модульности, тестируемости и повторного использования кода была спроектирована строгая многоуровневая архитектура. Система разделена на три уровня абстракции:

\begin{enumerate}
    \item Уровень арифметики конечного поля (Low-level Field Core): Базовый уровень, реализующий операции над элементами поля $GF(2^{255}-19)$ в представлении Radix-51. Здесь сосредоточены наиболее критичные к производительности операции: модульное умножение, возведение в квадрат и инверсия. Уровень поддерживает две реализации: скалярную и векторную (с использованием AVX2-интринсиков).
    
    \item Уровень групповых операций (Curve Geometry Layer): Средний уровень, оперирующий точками эллиптической кривой. Реализует формулы сложения и удвоения в проективных координатах, а также алгоритмы многократного сложения (скалярного умножения) точки. Данный уровень абстрагирован от конкретной реализации арифметики поля и использует унифицированные интерфейсы для работы как с одиночными точками, так и с пакетами точек (batching).
    
    \item Уровень криптографических примитивов (High-level Primitives): Верхний уровень, предоставляющий конечному пользователю готовые методы для реализации протоколов X25519 и Ed25519. На этом уровне реализуется логика, специфичная для конкретных алгоритмов: обработка закрытых ключей (Clamping), хеширование (SHA-512 для Ed25519), сериализация точек в байтовые массивы согласно RFC 7748/8032 и предоставление безопасного интерфейса на базе структур Span<byte>.
\end{enumerate}
Такое разделение позволяет независимо оптимизировать и тестировать базовую математику, не затрагивая высокоуровневую логику цифровой подписи.
\subsection{Проектирование структур данных вычислительного ядра}

Скорость криптографических вычислений во многом определяется тем, как элементы конечного поля представлены в памяти. Именно от структуры данных зависит, насколько эффективно будут загружены регистры процессора и кэш-память. 
В данной работе был применен подход «SIMD-first»: в отличие от традиционных библиотек, оперирующих одиночными числами, здесь базовым вычислительным блоком изначально является пакет из нескольких элементов. 
Это позволяет максимально задействовать возможности аппаратного ускорения с самого нижнего уровня иерархии вычислений \cite{DotNetPerformance}.


\subsubsection{Архитектура векторизованного представления Radix-51: AoS против SoA}

Векторизация криптографических алгоритмов сталкивается с фундаментальной проблемой организации данных в памяти. 
Традиционный подход, при котором каждый элемент поля (состоящий из пяти 64-битных лимбов Radix-51) является отдельным объектом, называется Array of Structures (AoS). В парадигме AoS четыре элемента поля в памяти располагаются последовательно:
$$[A_{L0}, A_{L1}, A_{L2}, A_{L3}, A_{L4}], [B_{L0}, B_{L1}, \dots] \dots $$

Попытка загрузить такие данные в 256-битный векторный регистр (вмещающий четыре 64-битных числа) приведет к тому, что в одном регистре окажутся разные лимбы одного и того же числа (например, $A_{L0}, A_{L1}, A_{L2}, A_{L3}$). Выполнять параллельные математические операции над таким вектором невозможно, так как алгоритм умножения требует умножения лимба $L_0$ одного числа строго на лимб $L_0$ другого.

Для решения этой проблемы в рамках проекта спроектирована архитектура данных на основе Structure of Arrays (SoA). Базовым строительным блоком вычислительного ядра является пользовательский тип значений FieldElement4, инкапсулирующий сразу четыре элемента конечного поля $GF(2^{255}-19)$.

Для достижения максимальной производительности в C\# структура спроектирована как непрерывный блок памяти фиксированного размера ([StructLayout(LayoutKind.Sequential)]). Вместо скалярных полей она содержит ровно пять векторных регистров типа Vector256<ulong> (лимбы $L_0, L_1, L_2, L_3, L_4$). Общий размер структуры составляет 160 байт, что гарантирует ее размещение исключительно на стеке вызовов (Stack) или непосредственно в регистрах процессора.

В парадигме SoA данные внутри регистров FieldElement4 располагаются следующим образом:
\begin{enumerate}
    \item Регистр $L_0$ содержит нулевые лимбы для четырех независимых чисел: $[A_{L0}, B_{L0}, C_{L0}, D_{L0}]$.
    \item Регистр $L_1$ содержит первые лимбы для четырех независимых чисел: $[A_{L1}, B_{L1}, C_{L1}, D_{L1}]$.
    \item И так далее до регистра $L_4$.
\end{enumerate}

Подобная архитектура позволяет вызывать базовые математические операции (сложение, вычитание, умножение) единожды, при этом процессор за один такт выполняет вычисление сразу для четырех криптографических контекстов с использованием интринсиков System.Runtime.Intrinsics.X86.Avx2.

\subsubsection{Паттерн проектирования для константного времени (Constant-Time SIMD)}

Архитектура слоя арифметики построена на принципах \textit{data-oblivious programming} (программирование, не зависящее от данных). В интерфейс базовой структуры не включены методы сравнения с ранним возвратом (early return).

Для операций, требующих условного выбора (например, условная замена точек в лестнице Монтгомери на основе бита секретного ключа), спроектирован векторизованный примитив CSwap. Вместо конструкции if, которая компилируется в инструкции условных переходов, логика реализована через аппаратные маски с использованием метода Avx2.BlendVariable.

Данная аппаратная инструкция принимает два сравниваемых вектора и вектор маски. Если бит маски равен $1$, берется байт из первого вектора, если $0$ — из второго. В методе CSwap вычисляется 256-битная маска swapBits (состоящая из участков 0x00...00 или 0xFF...FF в зависимости от того, требует ли конкретный из 4-х ключей обмена в данной итерации). 
Применение BlendVariable ко всем пяти регистрам структуры $L_0 \dots L_4$ позволяет выполнить независимый условный обмен четырех пар чисел абсолютно параллельно, без единого ветвления в коде. Это обеспечивает строгую математическую защиту от атак по сторонним каналам (Side-Channel Attacks).

\subsection{Взаимодействие с данными: транспонирование AoS и SoA}

Поскольку внешние ключи передаются пользователем в виде последовательных байтовых буферов (AoS), а вычислительное ядро оперирует векторами (SoA), необходим механизм их эффективного преобразования. В данной реализации логика транспонирования данных инкапсулирована непосредственно в методах FromBytes и ToBytes структуры FieldElement4.

При загрузке данных библиотека упаковывает четыре независимых 32-байтных значения в 256-битные регистры процессора. Это позволяет всем последующим алгоритмам (скалярному умножению, операциям на кривых Монтгомери и Эдвардса) работать исключительно с векторными типами данных (FieldElement4, PointExt4). По завершении вычислений результат автоматически сериализуется обратно в стандартный байтовый формат. Такой подход скрывает сложность векторной алгебры от верхних уровней библиотеки и минимизирует накладные расходы на подготовку данных.
\subsection{Модель управления памятью и паттерн Zero Allocation}

В экосистеме .NET процесс выделения памяти в управляемой куче (heap) посредством оператора new является быстрой операцией, однако последующая очистка памяти сборщиком мусора (GC) вызывает приостановку выполнения всех потоков приложения (Stop-The-World). В криптографических библиотеках старого поколения (например, Bouncy Castle) при сложении двух больших чисел каждый раз создается новый объект-массив (byte[] или uint[]), что при высоких нагрузках приводит к катастрофической деградации производительности (GC Churn).

Разработанная в рамках дипломной работы архитектура полностью опирается на парадигму Zero Allocation (нулевые аллокации). Для реализации этой парадигмы были применены следующие программные решения:

\subsubsection{Использование структур по ссылке (ref struct)}

Все базовые математические структуры (FieldElement4) спроектированы как типы значений. Для защиты от случайного копирования больших объемов данных при вызове методов повсеместно используются модификаторы передачи параметров по ссылке (ref и in). 
Это означает, что при вычислении $A \times B = C$ в методе умножения не создается новый объект $C$. Вместо этого функция принимает заранее выделенную на стеке область памяти для $C$ по ссылке и записывает результат напрямую туда:
static void Multiply(in FieldElement a, in FieldElement b, ref FieldElement result).

Криптографические ключи являются высокочувствительными данными. Архитектура библиотеки предусматривает строгий контроль жизненного цикла секретов. 

В процессе декодирования закрытого ключа (Private Key) и применения математической операции «зачистки битов» (Clamping, необходимой для устранения уязвимостей малых подгрупп), промежуточные скаляры хранятся исключительно в регистрах процессора и локальных переменных стека. После завершения функции генерации публичного ключа или подписи, механизмы среды CLR (Common Language Runtime) автоматически инвалидируют стек-фрейм, а библиотека гарантирует отсутствие "утечек" копий ключей в долгоживущие объекты.
